\section{Difesa e Validazione di
BlueAgent}\label{difesa-e-validazione-di-blueagent}

Come Sistema Operativo su cui basare lo sviluppo è stato scelto
\textbf{Debian 13}, in quanto offre un kernel stabile ed un'ottima
compatibilità tra i vari pacchetti e le componenti utilizzate da
BlueAgent, facilitando la configurazione di OVN, ZFS ed Incus.

Su questo sistema operativo si è scelto di utilizzare ZFS per creare lo
storage pool per le immagini dei container. ZFS offre il meccanismo CoW
a costo computazionale costante, ed è stato preferito rispetto ad
alternative come BTRFS poiché ritenuto notevolmente più maturo e
performante.

Si beneficia del CoW in due scenari chiave: quando si effettua il deploy
della challenge per un team (clonando l'immagine base dallo storage
pool) e durante la duplicazione per la verifica.

L'adozione di Incus è dovuta alla necessità di offrire ai concorrenti
macchine realistiche mediante system container (dotati di un proprio
init system completo e del processo \passthrough{\lstinline!init!} con
\passthrough{\lstinline!PID!} 1). Tale approccio permette ai
partecipanti di operare realisticamente accedendo a tutte le utility di
sistema e installando pacchetti software. Incus gestisce inoltre
nativamente gli storage pool ZFS e le reti OVN, semplificando
l'architettura del software BlueAgent.

Il principale vantaggio introdotto da BlueAgent è l'invisibilità totale
dell'Oracolo rispetto ai concorrenti: per i partecipanti è impossibile
intercettare o sniffare il traffico di rete generato durante la
validazione, poiché la rete di verifica \(\text{verify-}k\) risiede in
un dominio SDN logicamente separato. In
\passthrough{\lstinline!nftables!}, le regole di filtraggio consentono
il transito unicamente tra l'endpoint VPN \(\text{wg-}k\) e il bridge di
produzione \(\text{br-team-}k\).

Inoltre, l'adozione del meccanismo CoW riduce significativamente i tempi
di provisioning dei container come abbiamo visto in 5.4, mentre
l'impiego di OVN automatizza la configurazione della topologia di rete.
